% Ban thu: chi giu cac cau co trong CD2

%@ id=TT26.hadong-l1.003 | trang=1 | so=II | de=II | vitri=CD2.4,C3.9.5 | muc=8 | tt=0
\begin{bt}
Cho hai biểu thức $A = \dfrac{9-x}{\sqrt{x}+1}$ và $B = \dfrac{3\sqrt{x}-6}{x-9} - \dfrac{2}{\sqrt{x}+3}$ (với $x \ge 0, x \ne 9$).
1. Tính giá trị của biểu thức $A$ khi $x = 16$.
2. Chứng minh: $B = \dfrac{\sqrt{x}}{x-9}$.
3. Xét biểu thức $P = AB$. Tìm $x$ để $|P| - P = 0$.
\loigiai{
1. Thay $x = 16$ (TMĐK) vào biểu thức $A$ ta được:
\[ A = \dfrac{9-16}{\sqrt{16}+1} = \dfrac{-7}{5}. \]
Vậy $A = \dfrac{-7}{5}$ khi $x = 16$.
2. Ta có:
\[ B = \dfrac{3\sqrt{x}-6}{x-9} - \dfrac{2}{\sqrt{x}+3} = \dfrac{3\sqrt{x}-6-2\sqrt{x}+6}{x-9} = \dfrac{\sqrt{x}}{x-9} \quad \text{(đpcm)}. \]
3. Ta có:
\[ P = AB = \dfrac{9-x}{\sqrt{x}+1} \cdot \dfrac{\sqrt{x}}{x-9} = \dfrac{-\sqrt{x}}{\sqrt{x}+1}. \]
Vì $x \ge 0$ nên $\sqrt{x} \ge 0$ và $\sqrt{x}+1 > 0$, suy ra $P \le 0$.
Do đó $|P| - P = 0 \Leftrightarrow |P| = P \Leftrightarrow P \ge 0$.
Kết hợp với $P \le 0$ suy ra $P = 0 \Leftrightarrow \dfrac{-\sqrt{x}}{\sqrt{x}+1} = 0 \Leftrightarrow x = 0$ (TMĐK).
Vậy $x = 0$.
}
\end{bt}
